%%=============================================================================
%% Voorwoord
%%=============================================================================

\chapter*{\IfLanguageName{dutch}{Woord vooraf}{Preface}}%
\label{ch:voorwoord}

%% TODO:
%% Het voorwoord is het enige deel van de bachelorproef waar je vanuit je
%% eigen standpunt (``ik-vorm'') mag schrijven. Je kan hier bv. motiveren
%% waarom jij het onderwerp wil bespreken.
%% Vergeet ook niet te bedanken wie je geholpen/gesteund/... heeft

Deze bachelorproef kwam tot stand vanuit een concreet praktijkprobleem binnen E.S.C. BV: het veilig en beheersbaar migreren van klantdocumentatie van Microsoft OneNote naar Remote Desktop Manager. Tijdens dit traject werd duidelijk hoe groot de impact van documentatiestructuur is op dagelijkse IT-operaties. Wat op het eerste gezicht een technische migratieopdracht lijkt, blijkt in de praktijk ook een oefening in governance, toegangsbeheer en kwaliteitscontrole.

Voor mij was dit onderzoek bijzonder leerrijk omdat het theorie en praktijk nauw samenbracht. De literatuur rond Privileged Access Management, information overload en beveiligingsprincipes bood een inhoudelijk kader, terwijl de implementatie met PowerShell en Microsoft Graph API die principes concreet toetsbaar maakte. De uitdaging lag niet alleen in het bouwen van scripts, maar vooral in het uitwerken van een robuuste workflow die ook op grotere schaal bruikbaar blijft.

Graag wil ik mijn promotor en co-promotor bedanken voor hun begeleiding, kritische feedback en inhoudelijke ondersteuning tijdens dit traject. Daarnaast bedank ik mijn collega's bij E.S.C. BV voor de praktische input en het vertrouwen om deze case in een realistische context uit te werken. Hun ervaring uit het werkveld heeft mee vorm gegeven aan de keuzes in deze bachelorproef.

Tot slot hoop ik dat dit werk niet alleen een bruikbaar resultaat oplevert voor de huidige migratiecontext, maar ook een basis vormt voor verdere professionalisering van documentatiebeheer en kennisdeling binnen de organisatie.